\chapter{Current Status and Remaining Work}
\label{ch:status}

This chapter gives a clear snapshot of what is complete, what is prepared
but not yet executed, and what future work the evidence in this report
justifies. The three categories are kept distinct throughout: a workstream
is \emph{completed} only if it has produced the result it set out to
produce; \emph{prepared} if the groundwork is built and verified but the
next step has not been run; and \emph{recommended} if it is a justified
next step that has not yet been scheduled or started.

\section{Current Project Status}
\label{sec:status-current}

Table~\ref{tab:status-summary} summarises every workstream addressed in
this report.

\begin{table}[htbp]
\centering
\caption{Current status by workstream. "Completed" means the workstream
produced the result it set out to produce, including a negative or bounded
one; it does not mean every question in that area is closed.}
\label{tab:status-summary}
\footnotesize
\begin{tabular}{p{3.0cm}p{2.0cm}p{4.6cm}p{3.4cm}}
\toprule
Workstream & Status & Evidence/output & Next action \\
\midrule
Dataset audit/alignment & Completed & 791-row aligned basis, corruption-handling audit (Table~\ref{tab:corruption-handling}) & None required \\
Visible-corrosion regression & Completed & Three independent phases, campaign-level robustness confirmed (Section~\ref{sec:results-visible}) & None required \\
Structural-condition feasibility & Completed (feasibility study) & Mixed-input and metadata-only models, both regime-tested (Section~\ref{sec:results-structural}) & None required for this question as posed \\
Robustness diagnostics & Completed & Two independent phases, three regimes each (Section~\ref{sec:results-robustness}) & None required \\
Terminal ultimate-load refocus & Completed & Three-way ablation, subgroup and campaign checks (Section~\ref{sec:phase4b}) & None required \\
Degradation / proxy-RUL screening & Completed (methodological/diagnostic) & Two independently configured generations (Section~\ref{sec:results-degradation}) & Further methodological refinement possible with existing data (Section~\ref{sec:status-cleanup}); validated prognostic claims require additional longitudinal/failure data (Section~\ref{sec:status-structural-data}) \\
Four-class classification & \textbf{Prepared, not started} & Schema, split, leakage checks, augmentation complete; zero trained models (Section~\ref{sec:phase5}) & Train a baseline classifier (Section~\ref{sec:status-next}) \\
Software / reproducibility cleanup & Mixed & Some items done (feature QC, raw/smoothed trajectories); YAML bug and string corruption still live (Chapter~\ref{ch:engineering}) & See Section~\ref{sec:status-cleanup} \\
% INTERNAL TODO: this row is provisional and must be refreshed during the
% final, whole-report QA pass once Chapter 1, Chapter 10, the executive
% summary, abstract, and appendices are complete. Do not treat this row as
% frozen alongside the rest of this chapter.
Final report/documentation & In progress & Chapters 2--8 drafted and frozen; Chapter 9 (this chapter) provisional pending this row's refresh; Chapters 1, 10, executive summary, abstract, appendices not yet written & Continue per approved sequencing \\
\bottomrule
\end{tabular}
\end{table}

\section{Immediate Next Work}
\label{sec:status-next}

Among the research and modelling workstreams — as distinct from completing
this report itself, which remains the project's immediate task
(Table~\ref{tab:status-summary}) — the highest-priority next technical
step is training a baseline classifier for the four-class branch: it is
the only such workstream that is fully prepared and has produced no
model-performance evidence at all
(Section~\ref{sec:lim-classification}). The practice this report
recommends is already specified in Chapter~\ref{ch:framework} and is
restated here as a checklist, not invented: keep validation and test
specimen-disjoint, as already enforced by the saved split manifests
(Section~\ref{sec:eng-splits}); select any hyperparameter or
architecture choice using the validation partition only; evaluate the
test partition once, at the end, and report macro-averaged F1 and
per-class recall alongside a full confusion matrix rather than overall
accuracy alone, given the severity of the class imbalance
(Section~\ref{sec:metrics}); and apply class weighting or an equivalent
technique, since the existing augmentation does not itself rebalance the
classes (Section~\ref{sec:lim-classification}). A convolutional or
transformer-based architecture (for example a ResNet-50 or a Vision
Transformer) is a candidate choice consistent with the strategy already
named in Chapter~\ref{ch:framework}, not a committed or mandatory one; a
simpler baseline should establish a reference point before any more
complex architecture is attempted.

\section{Structural / Prognostic Work Required for Stronger Claims}
\label{sec:status-structural-data}

Model-engineering improvements and new data collection address different
gaps and should not be conflated. On the engineering side, exporting
genuinely out-of-fold structural predictions for the degradation stage,
persisting a complete feature contract with every trained model, and
checking the baseline structural estimates' physical plausibility are
concrete, currently unimplemented improvements
(Section~\ref{sec:eng-contracts}). The project has already run an
explicit image-only/metadata-only/combined feature-set ablation for
terminal ultimate load (Section~\ref{sec:phase4b}); the corresponding
ablation has not been repeated for the hidden-damage/degradation stage
specifically, but this report does not identify a distinct, specific
recommendation for it beyond the general out-of-fold and feature-contract
improvements just stated. None of these engineering items, however,
closes the evidence gap identified in Chapter~\ref{ch:interpretation}:
sparse, single-timepoint, campaign-confounded structural supervision
(Section~\ref{sec:lim-dataset}). Stronger structural or remaining-life
claims would require some combination of data this project does not have
— additional independent campaigns; design factors such as mesh
configuration and chloride level decoupled across campaigns rather than
perfectly aligned; more structural-labelled specimens; repeated or
intermediate structural measurements per specimen, rather than a single
terminal one; and an observed failure or service-life event time, if
remaining-useful-life is ever to be a directly supervised objective. These
are stated as requirements the current limitations point to, not as work
already planned or scheduled.

\section{Reproducibility Cleanup}
\label{sec:status-cleanup}

In order of concreteness rather than importance: quote the YAML
\texttt{treatment\_coarse: NO} value (or otherwise force it to load as a
string) to eliminate the boolean-parsing bug
(Section~\ref{sec:eng-schema}); locate and repair the "main\_first"
string-substitution pattern, prioritising the instances confirmed to break
execution (pandas aggregation calls, numeric YAML hyperparameters) over
the cosmetic ones (comments, the package version string); export an
explicit dependency specification for \texttt{main\_4}, matching the
practice already followed by \texttt{main}, \texttt{main\_2},
\texttt{main\_3}, and \texttt{augmentation}; persist a complete feature
list and hyperparameter record with every trained model, not only the
hidden-damage stage's existing partial list; and separate the remaining
analysis tables (for example \texttt{surface\_feature\_table.csv}) from a
model-ready feature matrix. None of these require new data collection;
all are confined to the existing codebase.

\section{Recommended Validation Before External Claims}
\label{sec:status-validation}

Before any claim from this project is used to support a decision outside
the project itself, this report recommends: validating any structural or
robustness conclusion against at least one additional, independent
campaign rather than the two available here; completing a strict,
single-pass test evaluation for the classification branch once a
classifier exists, with the test partition never consulted before that
point; regenerating degradation and proxy-RUL outputs from genuinely
out-of-fold structural predictions before interpreting those trajectories
as genuinely out-of-sample estimates — a separate requirement from, and
not a substitute for, the next one: explicitly justifying every proxy-RUL
threshold against an independent domain or physical criterion, since
out-of-fold structural inputs would improve the trajectories the
thresholds are applied to without themselves validating the threshold
values; repeating the subgroup and campaign-holdout
robustness checks already applied to terminal ultimate load for any new
target considered; and, where feasible, replicating key findings against
an external or independently collected dataset. This report does not
present the project, in its current state, as ready for deployment or for
generalisation beyond the two campaigns studied.
